% \iffalse
%<*driver>
\documentclass{l3doc}
\usepackage[T1]{fontenc}
\usepackage[brazilian]{babel}
\begin{document}
  \DocInput{abntex3.dtx}
\end{document}
%</driver>
% \fi
%
% \title{\pkg{abntex3}: infraestrutura mínima}
% \author{dantasrs}
% \date{2026-07-27 v0.0.0}
%
% \maketitle
%
% \begin{documentation}
%
% \section{Estado}
%
% Este arquivo produz o pacote mínimo usado para validar a cadeia de build,
% testes, instalação e distribuição do abnTeX3 Community. A versão 0.0.0 não
% implementa regras normativas e não oferece ainda uma API pública estável.
%
% \section{Uso}
%
% O pacote pode ser carregado sobre uma classe LaTeX usual.
%
% \begin{verbatim}
% \documentclass{article}
% \usepackage{abntex3}
% \begin{document}
% Documento mínimo.
% \end{document}
% \end{verbatim}
%
% A opção interna |debug| existe somente para exercitar o processamento de
% opções na infraestrutura de testes. Ela não integra a API pública e poderá
% desaparecer antes da primeira versão alfa.
%
% \end{documentation}
%
% \begin{implementation}
%
% \section{Implementação}
%
%    \begin{macrocode}
%<*package>
\NeedsTeXFormat{LaTeX2e}[2022-06-01]
\ProvidesPackage{abntex3}
  [2026-07-27 v0.0.0 Infraestrutura minima do abnTeX3 Community]

\newif\ifabntexthree@debug
\DeclareOption{debug}{\abntexthree@debugtrue}
\DeclareOption*{%
  \PackageWarning{abntex3}{Opcao desconhecida `\CurrentOption'}%
}
\ProcessOptions\relax

\ifabntexthree@debug
  \PackageInfo{abntex3}{Modo de diagnostico ativado}%
\fi
\endinput
%</package>
%    \end{macrocode}
%
% \end{implementation}
%
% \Finale
